% =====================================================================
%  config/style-fti.tex
%  Identité visuelle : couleurs de la charte graphique de l'Université
%  de Genève, en-têtes/pieds de page, mise en forme des titres.
%
%  Sources des couleurs (charte graphique UNIGE, Service de communication) :
%    https://www.unige.ch/communication/publier/charte/couleurs
%  Valeurs officielles :
%    - Rouge institutionnel UNIGE : Pantone 214C / #CF0063 / RVB 207,0,99
%    - Gris du sceau              : 45 % noir     / #A3A3A3 / RVB 163,163,163
%    - Couleur de la FACULTÉ de traduction et d'interprétation (FTI) :
%                                   Pantone 152C  / #FF5C00 / RVB 255,92,0
% =====================================================================

% ---------------------------------------------------------------------
%  Palette
% ---------------------------------------------------------------------
\definecolor{unigeRouge}{HTML}{CF0063} % rouge institutionnel (Pantone 214C)
\definecolor{unigeGris}{HTML}{A3A3A3}  % gris du sceau (45 % noir)
\definecolor{ftiOrange}{HTML}{FF5C00}  % couleur de la faculté FTI (Pantone 152C)

%  Couleur d'accent employée dans tout le document. Par défaut : la
%  couleur de la FTI (orange). Passez à \unigeRouge pour un rendu
%  « institutionnel » plus neutre.
\colorlet{accent}{ftiOrange}

% ---------------------------------------------------------------------
%  En-têtes et pieds de page (scrlayer-scrpage, KOMA)
% ---------------------------------------------------------------------
\usepackage[headsepline]{scrlayer-scrpage}
\pagestyle{scrheadings}
\clearpairofpagestyles

% Filet de séparation d'en-tête à la couleur d'accent.
\setkomafont{headsepline}{\color{accent}}
\KOMAoptions{headsepline=0.8pt}

% Contenu des en-têtes/pieds :
%   - en-tête extérieur : titre courant de chapitre/section
%   - pied extérieur    : numéro de page
\automark[section]{chapter}
\ihead{}
\ohead{\small\headmark}
\ofoot*{\normalfont\color{accent}\pagemark}
\cfoot*{}

% Marques de chapitre/section en style « propre » (sans « Chapitre N. »).
\renewcommand*{\chaptermarkformat}{}

% ---------------------------------------------------------------------
%  Mise en forme des titres (KOMA : \setkomafont / \addtokomafont)
% ---------------------------------------------------------------------
% Titres en police sans empattement, teintés à la couleur d'accent.
\addtokomafont{disposition}{\sffamily}
\addtokomafont{chapter}{\color{accent}}
\addtokomafont{section}{\color{accent}}
\addtokomafont{subsection}{\color{accent!85!black}}

% Numéro de chapitre imposant (style KOMA « chapterprefix » désactivé,
% numéro intégré au titre).
\renewcommand*{\chapterformat}{%
  \mbox{\textcolor{accent}{\Huge\thechapter}\enskip}%
}

% ---------------------------------------------------------------------
%  Liens hypertextes teintés à l'accent (discret, hyperref = hidelinks
%  masque les cadres ; ici on colore juste le texte des références).
% ---------------------------------------------------------------------
\hypersetup{
  colorlinks = false,   % pas de couleur agressive pour l'impression
  linkcolor  = accent,
  citecolor  = accent,
  urlcolor   = accent,
}

% ---------------------------------------------------------------------
%  Petites commandes utilitaires pour la page de titre / le style.
% ---------------------------------------------------------------------
% Filet horizontal à la couleur d'accent.
\newcommand{\ftirule}[1][\linewidth]{%
  {\color{accent}\rule{#1}{1.2pt}}%
}
